%%%
%
% $Autor: Wings $
% $Datum: 2021-05-14 $
% $Pfad: GitLab/CornerBLending $
% $Dateiname: Hints
% $Version: 4620 $
%
% !TeX spellcheck = de_DE/GB
%
%%%

\chapter{Hinweise}



Steps: 
\begin{itemize}
\item Was ist ein Gangbild?
\item Was soll gemessen werden  (was ist nützlich, um den Gang zu analysieren)
\item Welche Sensortypen werden verwendet? 
- Hüpfen laufen gehen muss man identifizieren können 
 Später: 

\item Verstehen, wie man Daten mit dem Mikrocontroller aufzeichnet. (Bsp. smart band) 
\end{itemize}
basics of the gait cycle and gate anaysis:  (\url {https://www.youtube.com/watch?v=1u6d1CX7o9c&ab_channel=Physiotutors})

- detailliert/ 
Explaining the Gait Cycle for the NPTE - 2021 (\url {https://www.youtube.com/watch?v=dvpi1WHCDwM&t=69s&ab_channel=ConqueringtheNPTE})


- Gangstörungen: Bewegungsstörung, die das Gehen bzw. das Gangbild betrifft. Sie kann neurologische, orthopädische oder psychogene Ursachen haben. Gangstörungen vermindern die Mobilität und erhöhen gerade bei älteren Patienten das Sturzrisiko.
%https://www.researchgate.net/publication/321515644_Gangstorungen


%Kurz beschrieben: 
%das Gangbild bezeichnet die visuelle Abbildung der Bewegung und Haltung der Extremitäten sowie deren Zusammenspiel mit dem Rumpf während des Gangzyklus bsw. des Gehens. 
%Bei einem gesunden Gang bewegen sich die Extremitäten fließend und koordiniert, während die Haltung des Oberkörpers ausbalanciert wird. was normalerweise zu einer nach vorne gerichteten Verlagerung des Schwerpunkts führt.
%unter ein Gangzyklus versteht man der Vorgang vom anfangs-Kontakt mit dem Boden bis zum nächsten anfangs-Kontak aus dem nächsten Zyklus. 
%ein bestimmte Gangbildmuster gibt es nicht denn jeder entwickelt seine eigene lösung : "Wie gelange ich mit einem Minimum an Anstrengung, mit ausreichend Stabilitat sowie einem guten Erscheinungsbild von einem Ort zum an- deren?"
%Dementsprechend, ist die Bestimmung der zu berücksichtigenden Parametern für den Rest des Projekts wichtig. %\cite{BeckersDeckers1997}



Sensoren
Gyroscope Winkel beschleunigung

Was ist das Problem?  
herausfoedwrungen: 
- generierung von data ( data aufzeichnen )
- beschränkte hardware : speicher, geh
Methode "KDD process" ( wenn man das Probelm kennt aber kein data hat )
 frage: wo soll ich die daten bearbeiten ( maschinelles lernen )
 rq: qualität sicherung bei VW wird mit bild verarbeitung gemacht 

1. Recherche: 
\begin{itemize}
\item \textbf{Part1} 
\item Was versteht man unter Gangbildern 
	 	- Using Machine Learning and Wearable Inertial Sensor Data for the Classification of Fractal Gait Patterns in Women and Men During Load Carriage \cite{U.Ahamed2021}
\item Wozu ist das Projekt 
		
		- Was bringt uns eine gangbildanalyse, wo wird die gebraucht 
\item Was ist der Zeit benutzt um Gangbildern zu bewerten ( tools , kosten, ort  …) 
\item \textbf{Part2} 
\item Arduino Nano BLE 33 definieren 
	\item Datasheet lesen
	\item Beschränkungen identifizieren : Speicher, Gehäuse, Komplexität, Geschwindigkeit.  
	\item Welcher Sensoren werden benutzt, warum und Welcher format von data ergeben die (IMU)  
\end{itemize}
2. Database erstellen
\begin{itemize}
\item intern : data generieren, 
\item extern:  
\end{itemize}
3. Data selection 
	 Nur Data von Hupfen Gehen und Laufen bleiben  
4. Data Preparation 
\begin{itemize}
\item qualitätprufung
\item Zeitbereich muss angepasst werden, gleiche länge ( chose the same time for each data )
\item Format anpassen  ( -> bild)
\end{itemize}
5. Data transformieren (frequenz analyse oder bild analyse) 


6. Data mining 
\begin{itemize}
\item Algorithm wählen (CNN)
\item Trenieren 
\item Evaluierung und Verifikation ( Verifikation Testen, Evaluierung: hab ich das richtige gemacht  )
\item Monitoring
\end{itemize}
7. Application 
\begin{itemize}
\item Beschreibung des Systems: handbuch ( extern), Entwicklung beschreibung (Teil der Arbeit)
\item Arduino HW 
\item Stromversorgung / Schrittstellen 
\item Software /IDE 
\item welche Daten ( allgemein, speziell)
 \item IMU /Datenstruktur 
\item Tests 
\item Musrererknennung / Auswertung (datenaufnahme/ anwendung des modeles/ Ergebnis)
\end{itemize}
rq 
- wie sollen meine daten aussehen ??
- was soll ich sammeln 


8.  Open questions /to do / conclusion 
Appendix 
Material lists 
Data sheet 
bibliography 
SW Ständer 
\begin{flushleft} 
	\includegraphics[width=15cm]{Images/GaitCycle}
\end{flushleft} 





